为目标任务选择源数据池，可能需要针对许多数据池组合分别训练和验证预测器。冻结特征允许更早比较候选，但保留好候选本身还不够：最终选择的质量与节省的工作必须能够说明预筛选的价值。本文评价一个两阶段协议，先用无标签源特征和目标特征形成短名单，再用源标签训练预测器，通过目标验证选出一个。我们比较目标加权实验设计（A-opt）、二阶矩匹配（MMD）和目标无关多样性。分析在共享线性模型下解释实验设计评分，并区分候选遗漏与验证误差。在固定样本成本的 DomainNet 组合上，评价50个 A-opt 或 MMD 候选的平均测试 Brier 低于评价227个随机候选。探索性地改变投影和读出器后，矩匹配仍有竞争力。固定参数重新分块后，域级匹配同样表现良好，而域内收益取决于构造和损失指标。PACS 和 Office-Home 的直接选择收益较小，可用改善空间也有限。独立计时表明，已有特征缓存时可以节省计算。加入实测特征提取成本后，逻辑回归仍有节省，岭回归的收益则发生反转。这些结果从最终选择质量和计算成本评价预筛选，而不只看候选保留率。
